\documentclass[12pt,a4paper]{article}

\usepackage[utf8]{inputenc}
\usepackage[T1]{fontenc}
\usepackage{lmodern}
\usepackage{xcolor}
\usepackage{geometry}
\usepackage{graphicx}
\usepackage{array}
\usepackage{booktabs}
\usepackage{hyperref}
\usepackage{amsmath}
\usepackage{setspace}
\usepackage{titlesec}
\usepackage{colortbl}

\geometry{margin=2.5cm}
\setstretch{1.15}

% Color theme
\definecolor{mainblue}{HTML}{1A4B84}
\definecolor{lightblue}{HTML}{E8F1FA}
\definecolor{darkgray}{HTML}{333333}

% Section styling
\titleformat{\section}
  {\Large\bfseries\color{mainblue}}
  {\thesection}{1em}{}

\titleformat{\subsection}
  {\large\bfseries\color{darkgray}}
  {\thesubsection}{1em}{}

% Hyperref styling
\hypersetup{
    colorlinks=true,
    linkcolor=mainblue,
    urlcolor=mainblue,
    citecolor=mainblue
}

% Title page
\makeatletter
\renewcommand{\maketitle}{
    \begin{center}
        \vspace*{   1cm}
        \colorbox{mainblue!15}{
            \parbox{\textwidth}{
                \centering
                {\huge\bfseries\color{mainblue} NoSQL Storage Cost Evaluation}\\[0.3cm]
                {\Large\color{darkgray} Full Automated Analysis and GUI Application}
            }
        }
        \vspace{1cm}

        {\large
        \textbf{Authors:}\\
        Myriam Ait Said \\
        Mohamed Aymane Sfouli \\
        Luca Rougemont \\
        George Shamieh
        
        }

        \vspace{0.7cm}
        {\normalsize ESILV        Big Data Structure}
        \vspace{1cm}

        \hrule height 0.8pt \vspace{0.4cm}
    \end{center}
}
\makeatother

\begin{document}

\maketitle

\section*{Introduction}

Dans le cadre du cours \emph{Big Data Structure}, nous avons étudié l’impact des différentes stratégies de dénormalisation NoSQL sur les coûts de stockage et la distribution lors du sharding.  
Au   delà des calculs demandés, nous avons choisi d'aller beaucoup plus loin : construire un outil complet, robuste et totalement générique capable d’analyser n’importe quel schéma NoSQL.

L’objectif initial était académique.  
Le résultat final est une application professionnelle, portable, élégante et utilisable par n’importe quel utilisateur non technique.

\bigskip


Notre démarche a été guidée par trois idées simples:
\begin{enumerate}
    \item Le stockage doit être mesuré avec précision, de manière transparente.
    \item La manipulation des schémas ne doit pas être réservée aux développeurs.
    \item L’outil doit rester générique, élégant, et simple à utiliser.
\end{enumerate}

C’est ainsi qu’est née notre application : \textbf{Big Data Structure     Storage Estimator}.

\section{Architecture du Projet}

Notre dossier final n’est pas une simple remise de scripts isolés.  
Il s’agit d’un \textbf{mini   projet logiciel structuré}, composé de plusieurs modules clairement définis.

\subsection*{1. AppCore/ — Le cœur de l’application}

\begin{itemize}
    \item \textbf{gui\_app.py} : l’interface graphique moderne réalisée en Tkinter.
    \item \textbf{stats.json} : les cardinalités globales (produits, stocks, orderlines...).
    \item \textbf{schema/} : tous les schémas utiles au projet (DB1 à DB5, exemple de la prof).
    \item \textbf{bigdata/models.py} : classes et abstractions décrivant les collections.
    \item \textbf{bigdata/sizes.py} : l’algorithme central de calcul de taille.
\end{itemize}

\subsection*{2. Lancement simplifié}

\begin{itemize}
    \item \textbf{Launch\_BigData\_Tool.vbs} : exécute l’application par un simple double   clic,
          sans console, sans terminal, sans manipulation.
\end{itemize}

\subsection*{3. Dossier Tests/}

Conçu pour valider la généricité de notre solution :

\begin{itemize}
    \item fichiers JSON de test structurés et non structurés,
    \item exemple fourni par l’enseignante,
    \item jeux de données personnalisés permettant d'inférer automatiquement un schéma.
\end{itemize}

L’application a été testée avec tous ces cas, chacun a parfaitement fonctionné.

\section{Fonctionnalités de l’Application}

Notre interface graphique offre une expérience fluide, professionnelle et autonome.

\subsection{1. Détection automatique des schémas}

Toute structure \texttt{\_schema.json} placée dans \texttt{AppCore/schema} apparaît instantanément dans la liste déroulante.

\subsection{2. Saisie intelligente du ``Document count''}

Le champ est rempli automatiquement en fonction du schéma :
\begin{itemize}
    \item DB1 \(\rightarrow\) nombre de produits,
    \item DB3 \(\rightarrow\) produits × entrepôts,
    \item DB4 \(\rightarrow\) nombre d’orderlines.
\end{itemize}

\subsection{3. Analyse complète du stockage}

L’outil :
\begin{itemize}
    \item lit et valide tout JSON Schema,
    \item calcule la taille exacte d’un document,
    \item convertit automatiquement en taille de collection et en gigaoctets.
\end{itemize}

\subsection{4. Inférence automatique d’un schéma}

Si l'utilisateur charge un \texttt{.json} qui n’est pas un JSON Schema,  
l’application propose automatiquement :

\begin{center}
\textit{``Ce fichier ne correspond pas à un schéma JSON. Voulez   vous inférer le schéma automatiquement ?''}
\end{center}

Cette fonctionnalité nous a permis de traiter l'exemple complet fourni par l’enseignante.

\section{Les Cinq Schémas Étudiés (DB1 à DB5)}

Chaque schéma représente une stratégie différente de dénormalisation :

\begin{itemize}
    \item \textbf{DB1 :} Product avec catégories + supplier.
    \item \textbf{DB2 :} DB1 + information de stock.
    \item \textbf{DB3 :} Stock comme racine, Product complètement embarqué.
    \item \textbf{DB4 :} OrderLine comme racine, avec Product complet.
    \item \textbf{DB5 :} Product embarque toutes les OrderLines liées.
\end{itemize}

Cette variété nous permet d'observer les avantages et limites de chaque modèle.

\section{Règles de Calcul du Stockage}

L’algorithme suit strictement les règles du cours :

\begin{itemize}
    \item Integer / Float : 8 bytes
    \item String : 80 bytes
    \item Date : 20 bytes
    \item LongString : 200 bytes
    \item Surcoût clé   valeur : 12 bytes
\end{itemize}

\subsection*{Approche algorithmique}

Le calcul est récursif :
\begin{enumerate}
    \item lecture du type,
    \item ajout du coût propre,
    \item traitement des objets internes,
    \item traitement des tableaux,
    \item propagation de la taille estimée,
    \item multiplication par la cardinalité totale.
\end{enumerate}

\section{Résultats du Calcul}

\begin{center}
\rowcolors{2}{lightblue!60}{white}
\begin{tabular}{lccc}
\toprule
\textbf{Database} & \textbf{Collection} & \textbf{Doc Size (bytes)} & \textbf{Collection Size (GB)} \\
\midrule
DB1 & Product & 1096 & 0.1021 \\
DB2 & Product & 1372 & 0.1278 \\
DB3 & Stock   & 1240 & 23.0968 \\
DB4 & OrderLine & 1464 & 5453.8250 \\
DB5 & Product & 1780 & 0.1658 \\
\bottomrule
\end{tabular}
\end{center}

\section{Interprétation}

\begin{itemize}
    \item \textbf{DB1 et DB2} : efficaces, équilibrées.
    \item \textbf{DB3} : duplication massive du Product.
    \item \textbf{DB4} : modèle non viable (plus de 5 To).
    \item \textbf{DB5} : charge maîtrisée malgré la dénormalisation.
\end{itemize}

\section{Sharding}

Nos scripts montrent que :

\begin{itemize}
    \item Un sharding par \#IDP est globalement excellent.
    \item \#IDW est très mauvais (fort skew).
    \item \#IDC distribue parfaitement les OrderLines.
\end{itemize}

\section{Validation avec l'Exemple Enseignante}

Le JSON de l’enseignante n’était pas un schéma.  
Grâce à notre inférence automatique, nous avons :

\begin{itemize}
    \item généré un schéma complet,
    \item calculé sa taille exacte,
    \item évalué son comportement au sein de notre pipeline.
\end{itemize}

Cette étape nous a permis de confirmer la généricité de l’application.

\section{Conclusion}
Pour cette première partie du projet, nous avons construit une applicatio portable et générale.  
Elle fonctionne avec n’importe quel schéma, n’importe quel exemple et n’importe quel volume de données.

Ce projet met en lumière :
\begin{itemize}
    \item notre compréhension des mécanismes de dénormalisation,
    \item notre maîtrise des problématiques de stockage,
    \item notre capacité à concevoir des outils fiables et professionnels.
\end{itemize}


\end{document}
